\documentclass[11pt]{article}

\usepackage{amsmath,amssymb,amsthm}
\usepackage{mathtools}

\newtheorem{theorem}{Theorem}
\newtheorem{remark}{Remark}

\newcommand{\OPT}{\mathrm{OPT}}
\newcommand{\WORST}{\mathrm{WORST}}

\begin{document}

\title{Bounds for \(P||\sum_j p_j C_j\) under adjacency restrictions}
\author{}
\date{}
\maketitle

\section{Preliminaries}

We consider the scheduling problem \(P||\sum_j p_j C_j\) on \(m\) identical parallel machines.  
Each job \(j\) has processing time \(p_j>0\), and the objective value of a schedule is

\[
\sum_j p_j C_j .
\]

Let \(P=\sum_j p_j\) and \(S=\sum_j p_j^2\). For a schedule whose machine loads are
\(L_1,\dots,L_m\), the objective value is independent of the order of the jobs on each
machine and is given by

\[
\sum_j p_j C_j
=
\frac12\left(\sum_{i=1}^m L_i^2+S\right).
\]

Indeed, if a machine contains jobs with processing times \(x_1,\dots,x_t\), then regardless
of the order,

\[
\sum_{r=1}^t x_r C_r
=
\sum_{r=1}^t x_r^2+\sum_{1\le r<s\le t}x_r x_s
=
\frac12\left(\left(\sum_{r=1}^t x_r\right)^2+\sum_{r=1}^t x_r^2\right).
\]

Let \(O_1,\dots,O_m\) denote the machines in an optimal schedule, and let

\[
L_i^{\OPT}=\sum_{j\in O_i}p_j
\]

be the load of OPT-machine \(i\). Let \(W_1,\dots,W_m\) denote the machines in the
considered List schedule, and let

\[
L_k^{\WORST}=\sum_{j\in W_k}p_j
\]

be the load of WORST-machine \(k\).

We assume throughout that jobs are labeled by their OPT-machine: job \(j\) has color \(i\)
if \(j\in O_i\). The adjacency restriction is:

\[
\text{If two jobs are consecutive on the same machine in WORST, then they have different OPT-colors.}
\]

We also assume the following job-size condition: there exists \(\beta>0\) such that for every
job \(j\in O_i\),

\[
p_j\le \beta \frac{L_i^{\OPT}}{m}.
\]

\section{The case of exactly \(m\) jobs per WORST machine}

\begin{theorem}
Suppose each WORST-machine contains exactly \(m\) jobs. Assume the adjacency restriction
and the job-size condition

\[
p_j\le \beta \frac{L_i^{\OPT}}{m}
\qquad
\text{for every } j\in O_i.
\]

Let

\[
h=\left\lceil \frac m2\right\rceil .
\]

Then

\[
\frac{\WORST}{\OPT}
\le
\frac{m\beta^2 h^2+1}{m+1}.
\]

Consequently,

\[
\frac{\WORST}{\OPT}
\le
\min\left\{
\frac{2m}{m+1},
\frac{m\beta^2 h^2+1}{m+1}
\right\},
\]

where \(2m/(m+1)\) is the general List-scheduling bound.
\end{theorem}

\begin{proof}
Fix a WORST-machine \(W_k\). Since \(W_k\) contains exactly \(m\) jobs, its sequence of
OPT-colors has length \(m\). Since no two adjacent jobs may have the same OPT-color, any
one color can appear at most

\[
h=\left\lceil \frac m2\right\rceil
\]

times in this sequence. Therefore,

\[
|O_i\cap W_k|\le h
\qquad
\text{for all } i,k.
\]

Using the job-size condition, we get

\[
L_k^{\WORST}
=
\sum_{i=1}^m \sum_{j\in O_i\cap W_k} p_j
\le
\sum_{i=1}^m |O_i\cap W_k| \cdot \beta \frac{L_i^{\OPT}}{m}.
\]

Since \(|O_i\cap W_k|\le h\), it follows that

\[
L_k^{\WORST}
\le
\frac{\beta h}{m}\sum_{i=1}^m L_i^{\OPT}
=
\frac{\beta hP}{m}.
\]

Hence

\[
\sum_{k=1}^m \left(L_k^{\WORST}\right)^2
\le
m\left(\frac{\beta hP}{m}\right)^2
=
\frac{\beta^2h^2P^2}{m}.
\]

On the other hand, by Cauchy--Schwarz,

\[
\sum_{i=1}^m \left(L_i^{\OPT}\right)^2
\ge
\frac{P^2}{m}.
\]

Since there are \(m\) machines and exactly \(m\) jobs on each WORST-machine, the total
number of jobs is \(m^2\). Therefore, again by Cauchy--Schwarz,

\[
S=\sum_j p_j^2
\ge
\frac{P^2}{m^2}.
\]

Now using the load formula for the objective,

\[
\frac{\WORST}{\OPT}
=
\frac{
\sum_{k=1}^m (L_k^{\WORST})^2+S
}{
\sum_{i=1}^m (L_i^{\OPT})^2+S
}.
\]

Substituting the upper bound on the WORST squared-load sum and the lower bound on the
OPT squared-load sum gives

\[
\frac{\WORST}{\OPT}
\le
\frac{
\frac{\beta^2h^2P^2}{m}+S
}{
\frac{P^2}{m}+S
}.
\]

Let \(X=P^2/m\). Then this becomes

\[
\frac{\WORST}{\OPT}
\le
\frac{\beta^2h^2 X+S}{X+S}.
\]

Assuming \(\beta h\ge 1\), which is the natural regime for this bound, the right-hand side
is decreasing in \(S\). Since \(S\ge P^2/m^2=X/m\), we obtain

\[
\frac{\WORST}{\OPT}
\le
\frac{\beta^2h^2 X+X/m}{X+X/m}
=
\frac{m\beta^2h^2+1}{m+1}.
\]

Thus

\[
\frac{\WORST}{\OPT}
\le
\frac{m\beta^2 h^2+1}{m+1}.
\]

Combining this with the general List-scheduling bound

\[
\frac{\WORST}{\OPT}\le \frac{2m}{m+1}
\]

gives the claimed minimum bound.
\end{proof}

\begin{remark}
The new adjacency-based bound improves the general List-scheduling bound precisely when

\[
\frac{m\beta^2 h^2+1}{m+1}
<
\frac{2m}{m+1}.
\]

Equivalently,

\[
m\beta^2h^2+1<2m,
\]

or

\[
\beta h < \sqrt{2-\frac1m}.
\]

Since \(h=\lceil m/2\rceil\), this condition is quite restrictive for \(m\ge 3\).
\end{remark}

\section{The case with between \(m\) and \(Q\) jobs per WORST machine}

We now consider the more general case where each WORST-machine contains at least \(m\)
jobs but at most \(Q\) jobs.

\begin{theorem}
Suppose each WORST-machine \(W_k\) contains \(q_k\) jobs, where

\[
m\le q_k\le Q.
\]

Assume the adjacency restriction and the job-size condition

\[
p_j\le \beta \frac{L_i^{\OPT}}{m}
\qquad
\text{for every } j\in O_i.
\]

Let

\[
H=\left\lceil \frac Q2\right\rceil .
\]

If the total number of jobs is \(n\), then

\[
\frac{\WORST}{\OPT}
\le
\frac{n\beta^2H^2+m}{n+m}.
\]

Consequently,

\[
\frac{\WORST}{\OPT}
\le
\min\left\{
\frac{2m}{m+1},
\frac{n\beta^2H^2+m}{n+m}
\right\}.
\]
\end{theorem}

\begin{proof}
Fix a WORST-machine \(W_k\). Its sequence has length \(q_k\). Since adjacent jobs cannot
have the same OPT-color, any one OPT-color can appear at most

\[
\left\lceil \frac{q_k}{2}\right\rceil
\]

times in \(W_k\). Since \(q_k\le Q\), we have

\[
|O_i\cap W_k|
\le
\left\lceil \frac{q_k}{2}\right\rceil
\le
\left\lceil \frac Q2\right\rceil
=
H.
\]

Using the job-size condition,

\[
L_k^{\WORST}
=
\sum_{i=1}^m \sum_{j\in O_i\cap W_k} p_j
\le
\sum_{i=1}^m |O_i\cap W_k| \cdot \beta \frac{L_i^{\OPT}}{m}.
\]

Since \(|O_i\cap W_k|\le H\), this gives

\[
L_k^{\WORST}
\le
\frac{\beta H}{m}\sum_{i=1}^m L_i^{\OPT}
=
\frac{\beta HP}{m}.
\]

Therefore,

\[
\sum_{k=1}^m \left(L_k^{\WORST}\right)^2
\le
m\left(\frac{\beta HP}{m}\right)^2
=
\frac{\beta^2H^2P^2}{m}.
\]

As before, by Cauchy--Schwarz,

\[
\sum_{i=1}^m \left(L_i^{\OPT}\right)^2
\ge
\frac{P^2}{m}.
\]

Also, since there are \(n\) jobs,

\[
S=\sum_j p_j^2
\ge
\frac{P^2}{n}.
\]

Using the load formula for \(P||\sum_j p_j C_j\), we obtain

\[
\frac{\WORST}{\OPT}
=
\frac{
\sum_{k=1}^m (L_k^{\WORST})^2+S
}{
\sum_{i=1}^m (L_i^{\OPT})^2+S
}
\le
\frac{
\frac{\beta^2H^2P^2}{m}+S
}{
\frac{P^2}{m}+S
}.
\]

Let \(X=P^2/m\). Then

\[
\frac{\WORST}{\OPT}
\le
\frac{\beta^2H^2X+S}{X+S}.
\]

Assuming \(\beta H\ge 1\), the right-hand side is decreasing in \(S\). Since

\[
S\ge \frac{P^2}{n}=\frac{mX}{n},
\]

we get

\[
\frac{\WORST}{\OPT}
\le
\frac{\beta^2H^2X+\frac{mX}{n}}{X+\frac{mX}{n}}
=
\frac{\beta^2H^2+\frac mn}{1+\frac mn}.
\]

Multiplying numerator and denominator by \(n\), this becomes

\[
\frac{\WORST}{\OPT}
\le
\frac{n\beta^2H^2+m}{n+m}.
\]

Thus

\[
\frac{\WORST}{\OPT}
\le
\frac{n\beta^2H^2+m}{n+m}.
\]

Combining with the general List-scheduling bound gives

\[
\frac{\WORST}{\OPT}
\le
\min\left\{
\frac{2m}{m+1},
\frac{n\beta^2H^2+m}{n+m}
\right\}.
\]
\end{proof}

\begin{remark}
The \(Q\)-bound improves the general List-scheduling bound precisely when

\[
\frac{n\beta^2H^2+m}{n+m}
<
\frac{2m}{m+1}.
\]

Equivalently,

\[
(m+1)(n\beta^2H^2+m)<2m(n+m).
\]

Rearranging gives

\[
n(m+1)\beta^2H^2
<
2mn+2m^2-m^2-m,
\]

hence

\[
n(m+1)\beta^2H^2
<
2mn+m^2-m.
\]

Therefore the \(Q\)-bound is better exactly when

\[
\beta H
<
\sqrt{
\frac{2mn+m^2-m}{n(m+1)}
}.
\]

Since

\[
H=\left\lceil \frac Q2\right\rceil,
\]

this becomes

\[
\beta \left\lceil \frac Q2\right\rceil
<
\sqrt{
\frac{2mn+m^2-m}{n(m+1)}
}.
\]
\end{remark}

\section{Specialization to \(n=m^2\)}

If \(n=m^2\), the \(Q\)-bound becomes

\[
\frac{\WORST}{\OPT}
\le
\frac{m^2\beta^2H^2+m}{m^2+m}
=
\frac{m\beta^2H^2+1}{m+1}.
\]

Thus, when \(n=m^2\),

\[
\frac{\WORST}{\OPT}
\le
\min\left\{
\frac{2m}{m+1},
\frac{m\beta^2H^2+1}{m+1}
\right\}.
\]

This recovers the earlier bound with \(H=h=\lceil m/2\rceil\) in the special case \(Q=m\).

\end{document}